% !TeX program = pdflatex
\documentclass[12pt,a4paper]{article}

\newcommand{\notelanguage}{finnish}
% Shared preamble for course notes and exercise sets.

\usepackage[T1]{fontenc}
\usepackage{lmodern}
\usepackage[english]{babel}

\usepackage[a4paper,margin=30mm]{geometry}
\usepackage{microtype}
\usepackage{graphicx}
\usepackage{enumitem}

\usepackage{amsmath}
\usepackage{amssymb}
\usepackage{amsthm}
\usepackage{mathtools}
\usepackage{aliascnt}

\usepackage[hidelinks,pdfusetitle]{hyperref}

\numberwithin{equation}{section}
\setlist{itemsep=0.25em,topsep=0.5em}

\theoremstyle{plain}
\newtheorem{theorem}{Theorem}[section]
\newaliascnt{lemma}{theorem}
\newtheorem{lemma}[lemma]{Lemma}
\aliascntresetthe{lemma}
\newaliascnt{proposition}{theorem}
\newtheorem{proposition}[proposition]{Proposition}
\aliascntresetthe{proposition}
\newaliascnt{corollary}{theorem}
\newtheorem{corollary}[corollary]{Corollary}
\aliascntresetthe{corollary}

\theoremstyle{definition}
\newaliascnt{definition}{theorem}
\newtheorem{definition}[definition]{Definition}
\aliascntresetthe{definition}
\newaliascnt{example}{theorem}
\newtheorem{example}[example]{Example}
\aliascntresetthe{example}
\newaliascnt{problem}{theorem}
\newtheorem{problem}[problem]{Problem}
\aliascntresetthe{problem}

\theoremstyle{remark}
\newaliascnt{remark}{theorem}
\newtheorem{remark}[remark]{Remark}
\aliascntresetthe{remark}

\usepackage[nameinlink,noabbrev]{cleveref}

\crefname{theorem}{theorem}{theorems}
\crefname{lemma}{lemma}{lemmas}
\crefname{proposition}{proposition}{propositions}
\crefname{corollary}{corollary}{corollaries}
\crefname{definition}{definition}{definitions}
\crefname{example}{example}{examples}
\crefname{problem}{problem}{problems}
\crefname{remark}{remark}{remarks}

\DeclareMathOperator{\im}{im}
\DeclareMathOperator{\rank}{rank}
\DeclarePairedDelimiter{\abs}{\lvert}{\rvert}
\DeclarePairedDelimiter{\norm}{\lVert}{\rVert}
\DeclarePairedDelimiter{\set}{\lbrace}{\rbrace}


\title{Johdatus logiikkaan I\\Harjoituksen 2 ratkaisut}
\author{}
\date{}

\begin{document}

\maketitle

\noindent

\begin{exercise}
  Oletataan, että \(v(p_0) = 1\), \(v(p_1) = 0\) ja \(v(p_2) = 1\). Määritä
\begin{enumerate}[label=\alph*), nosep]
  \item \(v((p_0 \lor \neg p_1) \rightarrow (p_1 \land p_2))\),
  \item \(v((p_0 \leftrightarrow \neg p_1) \land \neg p_2)\).
\end{enumerate}
\end{exercise}

\begin{solution}
\leavevmode\par
\begin{enumerate}[label=\alph*), nosep]
\item Totuusarvon määritelmän mukaan,
  \[
    \begin{aligned}
      v((p_0 \lor \neg p_1) \rightarrow (p_1 \land p_2)) &= v((1 \lor
      \neg 0) \leftarrow (0 \land 1)) \\
      &= v((1 \lor 1) \leftarrow 0) \\
      &= v(1 \leftrightarrow 0) \\
      &= 0.
    \end{aligned}
  \]

\item Totuusarvon määritelmän mukaan,
  \[
    \begin{aligned}
      v((p_0 \leftrightarrow \neg p_1) \land \neg p_2) &= v((1
      \leftrightarrow \neg 0) \land \neg 1) \\
      &= v((1 \leftrightarrow 1) \land 0) \\
      &= v(1 \land 0) \\
      &= 0.
    \end{aligned}
  \]
\end{enumerate}
\end{solution}

\begin{exercise}
Anna esimerkki totuusjakaumasta, jolla lause
\begin{enumerate}[label=\alph*), nosep]
\item \(((p_0 \rightarrow p_1) \land \neg(p_1 \rightarrow p_0))\)
\item \(((p_0 \land p_1) \leftrightarrow (p_0 \lor p_1))\)
\end{enumerate}
on tosi.
\end{exercise}

\begin{solution}
\leavevmode\par
\begin{enumerate}[label=\alph*), nosep]
\item Olkoon \(v\) totuusjakauma, jolla \(v(p_0) = 0\) ja \(v(p_1) = 1\).
Laskemalla lauseen totuusarvo, saadaan
\[
\begin{aligned}
  v((p_0 \rightarrow p_1) \land \neg(p_1 \rightarrow p_0)) &=
  v((0 \rightarrow 1) \land \neg(1 \rightarrow 0)) \\
  &= v(1 \land 1) \\
  &= 1.
\end{aligned}
\]

\item Olkoon \(v\) totuusjakauma, jolla \(v(p_0) = v(p_1) = 0\).
Laskemalla lauseen totuusarvo, saadaan
\[
\begin{aligned}
  v((p_0 \land p_1) \leftrightarrow (p_0 \lor p_1)) &=
  v((0 \land 0) \leftrightarrow (0 \lor 0)) \\
  &= v(0 \leftrightarrow 0) \\
  &= 1.
\end{aligned}
\]
\end{enumerate}
\end{solution}

\begin{exercise}
Onko lause
\begin{enumerate}[label=\alph*), nosep]
\item \(((p_0 \land p_1) \rightarrow (p_1 \rightarrow (p_2 \lor p_0)))\)
\item \((((p_0 \rightarrow (p_0 \rightarrow p_1)) \leftrightarrow (p_0
\rightarrow p_1))\)
\end{enumerate}
tautologia, kontingenssi vai ristiriita. Perustele.
\end{exercise}

\begin{solution}
\leavevmode\par
\begin{enumerate}[label=\alph*), nosep]
\item Ratkaistaan tehtävä totuustaulun avulla. Saadaan seuraava totuustaulu:
\[
\begin{array}{ccc|ccccccccc}
p_0 & p_1 & p_2
& (p_0 & \land & p_1)
& \boldsymbol{\to}
& (p_1 & \to & (p_2 & \lor & p_0)) \\
\hline
0 & 0 & 0
& 0 & 0 & 0
& \mathbf{1}
& 0 & 1 & 0 & 0 & 0 \\

0 & 0 & 1
& 0 & 0 & 0
& \mathbf{1}
& 0 & 1 & 1 & 1 & 0 \\

0 & 1 & 0
& 0 & 0 & 1
& \mathbf{1}
& 1 & 0 & 0 & 0 & 0 \\

0 & 1 & 1
& 0 & 0 & 1
& \mathbf{1}
& 1 & 1 & 1 & 1 & 0 \\

1 & 0 & 0
& 1 & 0 & 0
& \mathbf{1}
& 0 & 1 & 0 & 1 & 1 \\

1 & 0 & 1
& 1 & 0 & 0
& \mathbf{1}
& 0 & 1 & 1 & 1 & 1 \\

1 & 1 & 0
& 1 & 1 & 1
& \mathbf{1}
& 1 & 1 & 0 & 1 & 1 \\

1 & 1 & 1
& 1 & 1 & 1
& \mathbf{1}
& 1 & 1 & 1 & 1 & 1
\end{array}
\]
Nähdään että propositiolause on aina tosi, eli \emph{tautologia}.

\item Ratkaistaan tehtävä totuustaulun avulla. Saadaan seuraava totuustaulu:
\[
\begin{array}{cc|ccccccccc}
p_0 & p_1
& (p_0 & \to & (p_0 & \to & p_1))
& \boldsymbol{\leftrightarrow}
& (p_0 & \to & p_1) \\
\hline
0 & 0
& 0 & 1 & 0 & 1 & 0
& \mathbf{1}
& 0 & 1 & 0 \\

0 & 1
& 0 & 1 & 0 & 1 & 1
& \mathbf{1}
& 0 & 1 & 1 \\

1 & 0
& 1 & 0 & 1 & 0 & 0
& \mathbf{1}
& 1 & 0 & 0 \\

1 & 1
& 1 & 1 & 1 & 1 & 1
& \mathbf{1}
& 1 & 1 & 1
\end{array}
\]
Nähdään että propositiolause on aina tosi, eli \emph{tautologia}.
\end{enumerate}
\end{solution}

\begin{exercise}
Ovatko seuraavat lauseet loogisesti ekvivalentteja?
\begin{enumerate}[label=\alph*), nosep]
\item \((p_0 \rightarrow \neg p_1)\) ja \(\neg(p_1 \rightarrow p_0)\)
\item \((p_0 \lor p_1)\) ja \(((p_0 \rightarrow p_1) \rightarrow p_1)\)
\end{enumerate}
\end{exercise}

\begin{solution}
\leavevmode\par
\begin{enumerate}[label=\alph*), nosep]
\item Olkoon \(v(p_0) = v(p_1) = 0\) totuusjakauma. Laskemalla
propositiolauseille totuusjakaumat saadaan
\[
v(p_0 \rightarrow \neg p_1) = v(0 \rightarrow \neg 0) = v(0 \rightarrow 1) = 1,
\]
ja
\[
v(\neg(p_1 \rightarrow p_0)) = v(\neg(0 \rightarrow 0)) = v(\neg 1) = 0.
\]
Olemme löytäneet totuusjakauman \(v\), jolla propositiolauseiden totuusarvot
eriävät. Propositiolauseet eivät täten ole ekvivalentteja.

\item Piirretään propositiolauseiden totuustaulut. Saadaan
\[
\begin{array}{cc|ccc}
p_0 & p_1 & p_0 & \boldsymbol{\lor} & p_1 \\
\hline
0 & 0 & 0 & \mathbf{0} & 0 \\
0 & 1 & 0 & \mathbf{1} & 1 \\
1 & 0 & 1 & \mathbf{1} & 0 \\
1 & 1 & 1 & \mathbf{1} & 1
\end{array}
\]
ja
\[
\begin{array}{cc|ccccc}
p_0 & p_1 & (p_0 & \rightarrow & p_1) & \boldsymbol{\rightarrow} & p_1 \\
\hline
0 & 0 & 0 & 1 & 0 & \mathbf{0} & 0 \\
0 & 1 & 0 & 1 & 1 & \mathbf{1} & 1 \\
1 & 0 & 1 & 0 & 0 & \mathbf{1} & 0 \\
1 & 1 & 1 & 1 & 1 & \mathbf{1} & 1
\end{array}
\]
Totuustaulukoista lukemalla näemme propositiolauseiden totuusarvojen
oleman sama kaikilla mahdollisilla totuusjakaumilla, joten propositiolauseet
ovat \emph{ekvivalentteja}.
\end{enumerate}
\end{solution}

\begin{exercise}
Anna totuusjakauma, joka osoittaa allaolevan väitteen todeksi, tai todista,
että sellaista ei ole olemassa:
\begin{enumerate}[label=\alph*), nosep]
\item Propositiolause \((((p_0 \rightarrow \neg p_1) \land p_1)
\rightarrow p_0)\) on toteutuva.

\item Propositiolause \((\neg p_0 \rightarrow p_1) \lor \neg(p_0
\land \neg p_1)\) on kumoutuva.
\end{enumerate}
\end{exercise}

\begin{solution}
\leavevmode\par
\begin{enumerate}[label=\alph*), nosep]
\item Olkoon \(v(p_0) = v(p_1) = 0\) totuusjakauma. Laskemalla
propositiolauseen totuusarvo kyseisellä totuusjakaumalla saadaan
\[
\begin{aligned}
v(((p_0 \rightarrow \neg p_1) \land p_1)\rightarrow p_0)
&= v(((0 \rightarrow \neg 0) \land 0) \rightarrow 0) \\
&= v(((0 \rightarrow 1) \land 0) \rightarrow 0) \\
&= v((0 \land 0) \rightarrow 0) \\
&= v(0 \rightarrow 0) \\
&= 1.
\end{aligned}
\]
Propositiolause on täten \emph{toteutuva}.

\item Osoitetaan, ettei propositiolause ole \emph{kumoutuva}. Piirretään
propositiolauseelle totuustaulu:
\[
\begin{array}{cc|cccccccccc}
p_0 & p_1 & (\neg & p_0 \rightarrow & p_1) & \boldsymbol{\lor} & \neg
& (p_0 & \land & \neg & p_1) \\
\hline
0 & 0 & 1 & 0 & 0 & 0 & \mathbf{1} & 1 & 0 & 0 & 1 & 0 \\
0 & 1 & 1 & 0 & 1 & 1 & \mathbf{1} & 1 & 0 & 0 & 0 & 1 \\
1 & 0 & 0 & 1 & 1 & 0 & \mathbf{1} & 0 & 1 & 1 & 1 & 0 \\
1 & 1 & 0 & 1 & 1 & 1 & \mathbf{1} & 1 & 1 & 0 & 0 & 1
\end{array}
\]
Totuustaulusta nähdään, että propositiolause on tosi kaikilla totuusjakaumilla,
eli se ei ole \emph{kumoutuva}.
\end{enumerate}
\end{solution}

\begin{exercise}
Keksi kaksi propositiolausetta, joissa kummassakin esiintyy vähintään kaksi
eri propositiosymbolia ja jotka eivät ole loogisesti ekvivalentit, mutta
joista toinen on toisen looginen seuraus.
\end{exercise}

\begin{solution}
\leavevmode\par
\begin{enumerate}
\item Osoitetaan että propositiolauseet \((p_0 \leftrightarrow p_1)\) ja
\((p_0 \rightarrow p_1)\) täyttävät tehtävän vaatimukset. Ekvivalenssi on
tosi \emph{joss} \(v(p_0) = v(p_1)\). Näyssä tapauksissa myös implikaatio
on tosi. Toisaalta, kun \(v(p_0) = 0\) ja \(v(p_1) = 1\), ekvivalenssi
on epätosi, mutta implikaatio on tosi. Täten implikaatio on ekvivalenssin
looginen seuraus, mutta ekvivalenssi ja implikaatio eivät ole loogisesti
ekvivalentit.

\item Osoitetaan että propositiolauseet \((p_0 \land p_1)\) ja
\((p_0 \lor p_1)\) täyttävät tehtävän vaatimukset. Konjunktio on tosi
ainoastaan kun sekä \(p_0\) että \(p_1\) ovat tosia. Tässä tapauksessa
myös disjunktio on tosi. Toisin sanoen disjunktio on konjunktion looginen
seuraus. Toisaalta, \(v(1 \land 0) = 0\), mutta
\(v(1 \lor 0) = 1\). Täten disjunktio ja konjunktio eivät ole loogisesti
ekvivalentit.
\end{enumerate}
\end{solution}

\begin{exercise}
Oletetaan, että \(A\) ja \(B\) ovat propositiolauseita ja
\((A \rightarrow B) \implies (A \leftrightarrow B)\).
Näytä, että \(B \implies A\).
\end{exercise}

\begin{solution}
Oletuksen mukaan \(v(A \leftrightarrow B) = 1\) aina, kun
\(v(A \rightarrow B) = 1\). Mutta \(v(A \leftrightarrow B) = 1\) \emph{joss}
\(v(A) = v(B)\). Tästä voidaan päätellä, että
\[
v(A \rightarrow B) = 1 \implies v(A) = v(B).
\]
Olkoon \(v\) mielivaltainen totuusjakauma, jolla \(v(B) = 1\). Nyt
automaattisesti
\[
v(A \rightarrow B) = 1,
\]
koska implikaatio on epätosi joss \(v(A)\) on tosi ja \(v(B)\) epätosi.
Nyt oletuksen mukaan,
\[
v(A \leftrightarrow B) = 1,
\]
eli
\[
v(A) = v(B),
\]
eli \(v(A) = 1\). On siis osoitettu \(v(A) = 1\) aina, kun \(v(B) = 1\),
eli \(B \implies A\).
\end{solution}

\begin{exercise}
Olkoon \(v\) totuusjakauma, jolla \(v(p_i) = 1\) kaikilla \(i \in \NN\).
Olkoon \(A\) propositiolause, jossa ei esiinny negaatiota. Osoita, että
\(v(A) = 1\).
\end{exercise}

\begin{solution}
Todistetaan väite induktiolla rakenteen suhteen. Koska \(A\):ssa ei
esiinny negaatiota, riittää käsitellä propositiosymbolit ja
binäärikonnektiivit. Ensiksi, kun \(A = p_n\)
jollain \(n \in \NN\), saadaan \(v(A) = v(p_n) = 1\). Olkoon \(B\) ja \(C\)
propositiolauseita, jotka eivät sisällä negaatiota. Olkoon
\[
A = (B \circ C), \quad \circ \in \set{\land, \lor, \rightarrow,
\leftrightarrow}.
\]
Oletetaan induktio-oletuksena, että \(v(B) = v(C) = 1\).
Osoitetaan \(v(A) = 1\). Nyt
\[
v(B) = v(C) = 1 \implies v(B \circ C) = 1, \quad \text{kaikilla} \quad
\circ \in \set{\land, \lor, \rightarrow, \leftrightarrow},
\]
koska
\[
1 \land 1 = 1, \quad 1 \lor 1 = 1, \quad 1 \rightarrow 1 = 1, \quad 1
\leftrightarrow 1 = 1.
\]
Täten \(v(A) = 1\).

On osoitettu, että väite pitää paikkaansa propositiosymboleilla ja
säilyy totena kun rakennetaan suurempia lauseita binäärikonnektiivien avulla.
Näin ollen väite seuraa induktiosta rakenteen suhteen.
\end{solution}

\begin{exercise}
Mikä on totuusfunktio \(f_A\), kun
\[
A = ((p_0 \land \neg p_1) \lor \neg \neg(p_1 \rightarrow p_0))?
\]
\end{exercise}

\begin{solution}
Propositiolauseen \(A\) totuusfunktio \(f_A \colon \set{0, 1}^2
\rightarrow \set{0, 1}\) on
\[
\begin{aligned}
f_A(0, 0) &= f_A(1, 0) = f_A(1, 1) = 1 \\
f_A(0, 1) &= 0.
\end{aligned}
\]
\end{solution}

\begin{exercise}
Muodosta luvun 4 menetelmällä lause \(A\), jonka totuustaulu on
\[
\begin{array}{ccc|c}
p_0 & p_1 & p_2 & A \\
\hline
0 & 0 & 0 & 0 \\
0 & 0 & 1 & 1 \\
0 & 1 & 0 & 1 \\
0 & 1 & 1 & 0 \\
1 & 0 & 0 & 0 \\
1 & 0 & 1 & 1 \\
1 & 1 & 0 & 0 \\
1 & 1 & 1 & 1
\end{array}
\]
\end{exercise}

\begin{solution}
Luvun 4 Lemman (4.4) menetelmällä saadaan lause
\[
\begin{aligned}
A &= (\neg p_0 \land \neg p_1 \land p_2) \\
&\lor (\neg p_0 \land p_1 \land \neg p_2) \\
&\lor (p_0 \land \neg p_1 \land p_2) \\
&\lor (p_0 \land p_1 \land p_2).
\end{aligned}
\]
\end{solution}

\begin{exercise}[Haastetehtävä]
\end{exercise}

\begin{solution}
\end{solution}

\end{document}
